% Editable transcription; verify reused text against the original PDF.
\section{Modern AI techniques}

At the heart of the progress in AI lies the development of novel
techniques and algorithms commonly known as ``artificial neural
networks''.
In the late 20th century, ANNs started to become the model of choice
in many AI domains and are currently used in most modern AI applications.
Because of their importance in modern AI, we discuss their origin and
subsequent development, which shows a long history of mutual interaction
between and inspiration from psychology, neuroscience, and (artificial)
intelligence research.
Then, we discuss (artificial) reinforcement learning and other techniques
that are common in modern AI.

\subsection{Artificial neural networks}

Nowadays, many applications using artificial intelligence are powered,
at least partly, by (deep) artificial neural networks (ANNs).
ANNs come in many different varieties, but they all are composed of a
collection of largely identical processing ``units'' structured in a
hierarchical fashion, not unlike the way biological neurons make up
biological neural pathways.
In fact, psychology and neuroscience had, and still have, a strong
influence on the development of ANNs.

As their name suggests, the development of ANNs in the 1950s was
originally inspired by the structure and function of biological neurons
and networks (McCulloch \& Pitts, 1943).
The goal was to model cognition using models that stayed true to what,
at the time, was known about brain processing.
This approach, sometimes called ``connectionism'', is complementary
to the more symbolic approaches to artificial intelligence advocated
by, among others, the psychologists Alan Newell and Herbert Simon
(e.g. Newell \& Simon, 1956).

Initially, ANN models were composed of small arrays of artificial
neurons with a single input layer and one output layer, such as the
famous Perceptron model (Rosenblatt, 1958).
These models were often used to demonstrate that they can learn simple
logical functions (like the AND function) and embeddings (such as
recoding decimals to binary numbers).
Importantly, instead of explicitly programming the parameters
(or ``rules'') of the models (as was common in symbolic artificial
intelligence), the parameters (or ``weights'') of these networks
were learned.

After a period of reduced interest in and general pessimism about the
potential of artificial intelligence in general and ANNs in specific
(known as the first ``AI winter''), ANNs regained their popularity
when datasets and computing power increased throughout the first
decade of the 21st century.
Massive datasets, such as ImageNet (Deng et al. 2009), and the
technological developments that allowed networks to be trained on
powerful graphical cards allowed researchers to develop more complex
models, which often manifested as networks with many hidden layers
(i.e. deep neural networks; LeCun, Bengio, \& Hinton, 2015).
Also, a particular variant of ANNs, convolutional neural networks
(CNNs), were developed specifically to take advantage of locally
structured data (such as images and audio) by rectangular sets of
parameters, called filters, that were learned.
This made it unnecessary to carefully handcraft filters (as was common
in computer vision at the time).
CNNs proved to be extraordinarily powerful on visual tasks such as
object recognition, which was convincingly demonstrated when a CNN
beat the then state-of-the-art object recognition model during the
2012 ImageNet Large Scale Visual Recognition Challenge
(Krizhevsky, Sutskever \& Hinton, 2012).

Since 2012, research into ANNs has exploded, fueled both by a wider
interest from the AI community as well as the adaptation of ANNs
in different academic domains.
The general trend has been to develop more complex architectures
with more layers, often referred to as deep learning.
Extreme examples are the winning network from the ImageNet
competition in 2015, which featured over 150 hidden layers
(He, Zhang, Ren, \& Sun, 2016) and the aforementioned GPT-3 language
model featuring 96 layers and 175 billion parameters.
Apart from (computer) vision, ANNs have been adopted by other
scientific domains, including neuroscience
(Marblestone, Wayne, \& Kording, 2016), biology and medicine
(Ching et al. 2018), genomics (Zou et al. 2019), and geoscience
(Zhang, Zhang, \& Du, 2016).
Common to these applications is that (deep) neural networks are used
to estimate a (very) nonlinear mapping from a set of inputs to a set
of outputs.

The idea of stacking computations, as in the hierarchical layering
of (deep) ANNs, has been integrated in other AI domains such as
reinforcement learning and natural language processing (NLP).
This has given rise to deep reinforcement learning, used in the
previously mentioned AlphaZero model, and deep-learning-based NLP
models, such as the aforementioned GPT-3 model.
In these layered models, ANNs are used as a separate ``perceptual
module'' (which may or may not be pre-trained) that embeds information,
such as images or speech, in a way that a reinforcement learning
model or a language model can use the ANN as a basis.
Importantly, although reinforcement learning models may use ANNs
as part of their model architecture, they are fundamentally different
from ANNs in terms of goal.
Because of this, and because artificial reinforcement learning also
has roots in psychology and neuroscience, we now discuss this class
of models in more detail.

\subsection{Reinforcement learning}

Artificial reinforcement learning models learn optimal behavior
in a given context, through experience.
In contrast to ANNs, which tend to excel at perceptual tasks,
reinforcement learning is used to make an artificial agent learn
to behave.
Unsurprisingly, much of the progress in artificial reinforcement
learning was inspired by and used concepts from animal and human
learning (Sutton \& Barto, 1998, see chapter 11).
Early artificial reinforcement learning models, which did not feature
ANN modules, directly adapted or extended models from animal learning.
For example, the extended Rescorla-Wagner classical conditional model
(Rescorla \& Wagner, 1972), called temporal difference learning
(Sutton \& Barto, 1987), is extensively used in both psychology
(e.g. Gershman, Pesaran, \& Daw, 2009), neuroscience
(Neftci \& Averbeck, 2019) and AI research
(e.g. deep Q learning; Mnih et al. 2015).

These early reinforcement learning models were often evaluated
on relatively simple problems using a relatively simple environment
(i.e. with few ``states'') and a restricted and discrete set of possible
behaviors (i.e. with few ``actions'').
Examples of such ``toy problems'' are navigating a maze (in animal
learning) or learning to choose the most rewarding (virtual) slot
machine (in artificial reinforcement learning).
Driven by the desire to scale up reinforcement learning models
to complex environments and behaviors, 21st century researchers
started to add ANNs to their reinforcement learning models.
In these models, ANNs function as a preprocessing step that embeds
a complex (possibly continuous) environment as a more tractable set
of representations that can be fed into the reinforcement learning
model.
By adding ANNs researchers could scale up reinforcement learning
models that in turn led to the development of systems capable of
(super) human performance on a range of tasks, including playing
board games (Silver et al. 2018) and video games (Mnih et al. 2015),
as well as autonomous driving (Kiran et al. 2020).

\subsection{Other techniques}

ANNs and reinforcement learning are arguably the dominant techniques
used in modern AI models, but far from the only ones.
Here, we briefly discuss two other important techniques:
generative modelling and Bayesian models.

Within artificial intelligence, generative modelling is a broad set
of models that are in fact (often) a specific implementation of a
neural network.
The difference with ``standard'' neural networks is that generative
neural networks try to learn the distribution of the data (X) given
a particular feature (y), p(X \textbar{} y), in contrast to standard
(or ``discriminative'') neural networks that try to predict a
particular feature given the data, p(y \textbar{} X).
For example, instead of predicting object category (y) from images
(X), a corresponding generative model could try to learn what a
particular image (X) of a particular object (y) would look like.
For example, researchers trained a generative model on images of
faces using a generative adversarial network
(GAN; e.g. Karras, Aila, Laine, \& Lehtinen, 2017) and were then able
to generate photorealistic images of faces of non-existent people
(convincingly demonstrated on
\url{https://thispersondoesnotexist.com}).
Generative models have also been successfully applied in NLP,
with the recent GPT-3 model as a striking example
(Brown et al. 2020), and across domains, such as the text-to-image
model DALL-E, which is able to create images from text descriptions
(see Fig. 1).

Bayesian modelling (or Bayesian reasoning) is another technique
worth mentioning.
Although Bayesian techniques are used more in psychology than in AI,
many have argued that Bayesian modelling is crucial to overcoming
the current limitations of AI
(Lake, Ullman, Tenenbaum, \& Gershman, 2017;
Marcus, 2018; Pearl, 2019).
What's more, the abstract problem domains that Bayesian models are
often applied to are much closer to the problems presented in human
intelligence tests than those presented to ANNs.

Bayesian models come in many forms, but their common denominator
is that they aim to formalize interpretable and generalizable
mechanisms and knowledge.
In contrast to most ANNs, which are trained on a very specific task
(e.g. object recognition) and often need massive amounts of data,
Bayesian models are designed to learn generalizable mechanisms
that need little data to train
(Tenenbaum, Kemp, Griffiths, \& Goodman, 2011).
In a way, this approach resembles symbolic AI, with the major
difference being that modern Bayesian techniques are able to handle
more complex data and models.
To achieve models that can better generalize across domains
(``far transfer''), Bayesian models stress the importance of
abstract knowledge, functioning as priors in the Bayesian framework,
that constrain their structure and learning process.
A well-known application of Bayesian modelling is in the realm
of intuitive physics
(e.g. Battaglia, Hamrick, \& Tenenbaum, 2013).
Here, Bayesian models are used to model how humans (learn to)
understand the physical world and its physical laws in a causal way.
While Bayesian modelling is still mostly confined to realm of
psychology, many have advocated integrating ``connectionism''
(represented by ANN modelling) and ``symbolic AI'' (represented by
symbolic/Bayesian models; Lake et al. 2017; Marcus, 2018), and initial
efforts have produced some promising results
(Mao, Gan, Kohli, Tenenbaum, \& Wu, 2019; Yi et al. 2018).
The field of intelligence also sees a renewed interest in cognitive
modelling (e.g. Frischkorn \& Schubert, 2018), which may provide
new links between AI and intelligence research.

To summarize, modern AI relies on a set of techniques that are able
to solve complex tasks.
ANNs are extraordinarily effective in the context of perceptual
tasks (such as object and speech recognition), but are increasingly
used as modules in other applications as well, for example by
combining them with reinforcement learning and NLP models.
Whereas ANNs are particularly suited to model perceptual processes,
reinforcement learning represents a set of techniques that are very
effective in terms of learning adaptive behavior and can outperform
humans at a variety of tasks (including board and video games).
In addition, generative modelling complements techniques for
prediction by generating data, and Bayesian reasoning may provide
AI with models that alleviate AI of its current shortcomings such
as limited generalizability and data-intensive training regimes.
